\pdfoutput=1
\documentclass[11pt]{article}
\usepackage[utf8]{inputenc}
\usepackage[T1]{fontenc}
\usepackage{lmodern}
\usepackage[margin=0.95in]{geometry}
\usepackage{amsmath,amssymb,amsthm,mathtools}
\usepackage{microtype}
\linespread{0.98}
\addtolength{\textheight}{0.16in}
\usepackage[hidelinks]{hyperref}
\newtheorem{theorem}{Theorem}
\newtheorem{corollary}[theorem]{Corollary}
\hypersetup{pdftitle={Semiclassical Saturation of the Virasoro Vacuum Whittaker Ray},pdfauthor={}}
\title{Semiclassical Saturation of the Virasoro Vacuum Whittaker Ray}
\author{}
\date{}
\begin{document}
\maketitle
\begin{abstract}
The Virasoro vacuum Whittaker ray with \(L_2\)-eigenvalue \(r\) has a single quadrupole in its stress expectation and is strictly outside the classical vacuum orbit at finite central charge. We prove that this gap is subextensive on every macroscopic ray \(r=c\alpha\): the normalized energy converges to the lowest Mathieu eigenvalue formula for the classical vacuum orbit, for arbitrary real \(c\to\infty\) and every \(\alpha>0\). The normalized logarithmic norm is the integral of the resulting energy, and all fixed smooth normalized stress fluctuations vanish. At integer charge the fixed-copy reduced states in the current realization converge in trace norm to tensor powers of the corresponding smooth welded vacuum ray. The proof combines diagonal vacuum projection and real-charge subadditivity with a level-four current identity preventing energy escape, and a finite tensor certificate forcing the exterior Grunsky relations on de Finetti components. No classical irregular-block exponentiation is required.
\end{abstract}

\section{Statement and scope}

We use \( [L_m,L_n]=(m-n)L_{m+n}+c(m^3-m)\delta_{m,-n}/12\), \(L_n^*=L_{-n}\), and the irreducible unitary zero-weight vacuum space with \(L_{-1}\Omega=0\). Let \(\Psi_{c,r}\) have vacuum coefficient one and be annihilated by all positive modes except \(L_2\Psi_{c,r}=r\Psi_{c,r}\). Expectations are in its normalized ray. Thus the cylinder stress has the exact form
\[
 t_{c,r}(\theta)=\langle L_0\rangle_{c,r}-c/24+2r\cos(2\theta).
\]
A shifted \(L_2\)-phase amounts to a rotation. Put \(Z_c(r)=\|\Psi_{c,r}\|^2\), and for \(q\in\mathbb R\) let \(a_*(q)\) be the lowest eigenvalue of \(-\partial_\theta^2+2q\cos(2\theta)\) on \([0,\pi]\) with both function and derivative multiplied by \(i\) across the endpoints.

\begin{theorem}\label{thm:main}
For every \(\alpha>0\) set \(e_*(\alpha)=(1/4-a_*(6\alpha))/6\). With real \(c\ge2\), as \(c\to\infty\),
\begin{align}
 \frac{\langle L_0\rangle_{c,c\alpha}}{c}&\longrightarrow e_*(\alpha),\label{eq:energylimit}\\
 \frac{\log Z_c(c\alpha)}{c}&\longrightarrow
                  \int_0^\alpha e_*(u)\,\frac{du}{u}.\label{eq:freelimit}
\end{align}
The spectral measurement \(L_0/c\) concentrates exponentially about \(e_*(\alpha)\) at each fixed positive tolerance. The variance of \(c^{-1}T(g)\) converges to zero for each fixed smooth real circle test \(g\), where \(T(g)\) is the smeared cylinder stress. At any fixed \(c\ge2\) and \(r>0\) one has \(\langle L_0\rangle_{c,r}>c\,e_*(r/c)\).
\end{theorem}

\begin{corollary}\label{cor:product}
Realize integer charge \(N\) in \(N\) independent zero-momentum chiral currents. In the normalized ray \(\Psi_{N,N\alpha}\), the reduced state of each fixed \(k\) current copies converges in trace norm to \(P_\alpha^{\otimes k}\). Here \(P_\alpha\) is the one-current smooth diffeomorphism-vacuum ray with cylinder stress \(-1/24+e_*(\alpha)+2\alpha\cos(2\theta)\).
\end{corollary}

The finite-charge exclusion and sharp classical profile were established in \cite{p34}. The result here says that the strict obstruction disappears in the normalized macroscopic energy while remaining exact at every finite \(c\), and identifies both the limiting fluctuations and the current-copy hierarchy. It neither gives a finite-size correction nor shows that the full \(N\)-copy vector approaches a vacuum orbit in norm. In the Brown--Henneaux interpretation \cite{bh,cotler}, the result concerns one chiral vacuum module and its boundary stress, not a claim of a specific bulk path integral or a theorem for all CFT states.

Classical irregular conformal blocks already yield important Mathieu relations \cite{piatek14,piatek15}. The usual rank-one irregular exponential limit is not assumed here. In particular the Virasoro even-mode subalgebra generated by \(\widehat L_n=\tfrac12 L_{2n}+(c/16)\delta_{n0}\) has central charge \(2c\) and highest weight \(c/16\). Projecting to its cyclic highest-weight subspace yields a rank-one ray with parameter \(r/2\), but does not determine the full diagonal vacuum norm: the full vacuum quotient contains additional odd-mode descendants from level six. Fixed-parameter Nekrasov convergence \cite{nekrasov} is not the simultaneous limit of this vacuum quotient.
\section{\texorpdfstring{The macroscopic limit exists at all real \(c\)}{The macroscopic limit exists at all real c}}\label{the-macroscopic-limit-exists-at-all-real-c}

The vacuum Whittaker vectors are characterized by vacuum coefficient one and \[
 L_2\Psi_{c,r}=r\Psi_{c,r},\quad L_n\Psi_{c,r}=0\quad (n>0,n\ne2),\quad
 Z_c(r)=\|\Psi_{c,r}\|^2.
\]

To make the domain and normalization explicit, in the zero-charge \(c=1\) current representation take the oscillator vector
\[
 \Xi_r=\sum_{j\ge0}\frac{(2r)^j}{(2j)!}a_{-1}^{2j}\Omega,\qquad
 \|\Xi_r\|^2=\cosh(2r).
\]
Among the positive Sugawara modes only \(L_2=\tfrac12a_1^2\) acts nontrivially on this one-mode sector, with \(\Xi_r\) as eigenvector of eigenvalue \(r\). Tensor with the unitary vacuum of charge \(c-1\ge1\), then project orthogonally to the cyclic diagonal Virasoro vacuum module. This projection intertwines the Virasoro modes on finite-energy vectors and on the appropriate graph domains by the usual stress energy bounds \cite{carpi,fewster}. Every power of the energy has finite norm on \(\Xi_r\). Descendant pairings against all negative-mode monomials give uniqueness; the grade \(2j\) vector is proportional to \(r^j\) and the odd grades vanish. This recalls the construction of \cite{p34} inside the actual irreducible vacuum module, not a replacement Verma module.

We take \(r\ge0\) real; rotation handles a complex phase. Put \(f_c(s)=c^{-1}\log Z_c(ce^s)\). The Whittaker diagonal projection of \(\Psi_{c,r}\otimes\Psi_{d,u}\) is exactly \(\Psi_{c+d,r+u}\), by the descended mode equations and unit vacuum coefficient. Hence \(c f_c(s)\) is subadditive for any real \(c,d\ge2\). At fixed \(s\), \(0\le f_c(s)\le2e^s\) by the current realization above. Write a large real charge \(c=k c_0+\ell\) for any fixed \(c_0\ge2\), with \(\ell\in[2,2+c_0)\). The upper bounded remainder proves the continuous Fekete assertion \[
 f(s)=\lim_{c\to\infty,\ c\in\mathbb R} f_c(s)=\inf_{c\ge2}f_c(s).
\] The \(f_c\)'s are convex in \(s\), since each orthogonal grade enters with a nonnegative coefficient and \(r^{2j}\) on energy \(2j\). Their derivative is \(f_c'(s)=\langle L_0\rangle_{c,ce^s}/c\). The exact moment generating function of \(L_0\) is \(Z_c(re^t)/Z_c(r)\). On a fixed compact \(s\)-interval any unit-length exponential tilt and \(0\le f_c\le2e^s\) give uniform \(O(c^k)\) moments of \(L_0^k\) for every fixed \(k\); no differentiability of \(f\) is required for this bound.

\section{Current-site compactness and identification of stress expectations}\label{current-site-compactness-and-identification-of-stress-expectations}

For integer \(c=N\), realize the diagonal vacuum in \(N\) identical zero-momentum current Fock spaces, \(a_{n,i}^\dagger=a_{-n,i}\), and use the normalized \(\Psi_{N,N e^s}\). Write \(b_{n,i}=a_{n,i}/\sqrt n\) on \(n>0\). The state is invariant under orthogonal transformations of the copy index, because the stress modes are the sum of the Sugawara modes and the product vacuum is invariant. Its first \(k\) site reduced state is supported on the cyclic \textbf{diagonal} Virasoro vacuum subspace of those sites: each global generating mode is the sum of the corresponding local-diagonal and complementary-diagonal modes, and the two sectors commute. This support statement is at finite \(N\) and passes in trace norm for fixed \(k\).

The bound on \(L_0/N\) gives compactness of fixed-site reduced states in trace norm (each local oscillator-energy cutoff is finite-dimensional). Choose a consistent subsequential infinite hierarchy and use the strong bosonic quantum de Finetti representation \cite{definetti} \(\rho_k=\int\sigma^{\otimes k}\,d\mu(\sigma)\). One may obtain it by the finite-dimensional theorem on increasing finite local energy projectors with the above tightness; the normal infinite-system version is equivalent. Trace norm alone is \textbf{not} enough to pass the unbounded \(L_0\) or \(L_n\). Here the following independent estimates are decisive.

\textbf{No frequency escape.} In \(N\) copies set \(X=\sum_i (a_{-2,i})^2\Omega\). In the diagonal vacuum its level-four state projection is \(A L_{-4}\Omega+B L_{-2}^2\Omega\), where \(A=4(N+8)/(5N+22)\), \(B=-24/(5N+22)\). Indeed the overlaps with the two basis vectors are \(4N,0\) and their Gram entries are \(5N,3N,N(N+8)/2\). The orthogonal remainder has vanishing diagonal vacuum matrix elements for all of its field modes: the current Fock field for each Virasoro descendant in the vacuum subalgebra obeys the ordinary vacuum vertex algebra, and acting or taking the adjoint of its mode preserves that subspace. Concretely, for \(v,w\) in the diagonal subalgebra use \(Y(X,z)v=e^{zL_{-1}}Y(v,-z)X\). After taking a coefficient in the matrix element against \(w\), move the mode of \(Y(v,-z)\) using the unitary adjoint formula. The resulting vector paired against the orthogonal remainder is in the diagonal subalgebra \cite{voa}. No primary-state assumption on the remainder is needed. The energy-preserving zero mode identities are \[
 X_0=-2\sum_{i,m\ge1}(m^2-1)a_{-m,i}a_{m,i},\quad
 (L_{-4}\Omega)_0=3L_0,\quad
 (L_{-2}^2\Omega)_0=L_0^2+2L_0+2\sum_{n\ge1}L_{-n}L_n .
\] Their derivation uses the usual cylindrical degree-zero mode of the weight-four plane fields, equivalently \(z^4 X(z)\) constant Fourier coefficient; for \(X(z)=\sum_i:(\partial J_i(z))^2:\) each paired \(m,-m\) gives \(-2(m^2-1)\). For the other two fields \(Y(L_{-4}\Omega,z)=\frac12\partial^2 T(z)\) and \(Y(L_{-2}^2\Omega,z)=:T(z)T(z):\). Normal ordering at \(-2\) gives the displayed \(2L_0\) from \(L_1L_{-1}\). First apply on finite total-energy sectors, where only finitely many oscillator labels occur. The Whittaker vector's entire graded moment bounds and \(\sum_{n>0}\|L_n\Psi\|^2=r^2\|\Psi\|^2\) pass the expectation to the full vector by energy truncation. Thus \[
 2\sum_{i,m}(m^2-1)\langle a_{-m,i}a_{m,i}\rangle
   =\frac{24}{5N+22}(\langle L_0^2\rangle+2\langle L_0\rangle+2r^2)
       -\frac{12(N+8)}{5N+22}\langle L_0\rangle=O(N).
\] This gives uniform high-frequency tail in the one-site energy, \(O(K^{-2})\) after frequency \(K\), a condition not implied by the total-energy bound alone.

\textbf{No fixed-frequency occupation escape.} At a fixed frequency \(p\), let \(N_p=\sum_i b_{p,i}^\dagger b_{p,i}\) and \(Q_p=\sum_i b_{p,i}^2\). Averaging the one-site factorial second moment over the real species unit sphere by \(O(N)\) invariance gives \[
 \langle b_{p,1}^{\dagger\,2}b_{p,1}^{2}\rangle
    =\frac{2\langle N_p(N_p-1)\rangle+\langle Q_p^\dagger Q_p\rangle}{N(N+2)}.
\] The operator \(Q_p^\dagger Q_p\) on the \(N_p=n\) sector is bounded by \(n(n+N-2)\). To verify, use \([Q_p,Q_p^\dagger]=4N_p+2N\) and recursively compare the largest eigenvalue of \(Q_p^\dagger Q_p\) at \(n\) with that of \(Q_pQ_p^\dagger\) at \(n-2\) (the two have the same nonzero singular-value squares). Start at \(n=0,1\). Since \(pN_p\le L_0\), the factorial one-site moments are uniformly bounded. For any fixed finite set of modes this bounds the second moment of its local occupation sum; local quadratic fields on fixed modes therefore have uniform integrability. The tail of \(\langle a^\dagger_p a_{p+n}\rangle\) in a fixed stress mode can then be bounded by Cauchy--Schwarz using the high-frequency energy above. The anomalous part of a fixed stress Fourier mode only pairs finitely many positive indices of fixed sum. The same argument works in joint expectations with a bounded observable on disjoint sites, since it commutes with that local mode and contributes only its norm. In particular one-site means of each \(L_n\), including \(L_0\), pass to the de Finetti hierarchy.

\textbf{Every limiting state has the same nonzero stress modes.} Let \(B\) be finite rank with finite-energy range and adjoint range on the first \(k\) site spaces, with bounded norm. For \(n>0\), on the global ket \((L_n^{\rm tot}-\delta_{n2}r)\Psi=0\), and the corresponding adjoint fluctuation on the ket has squared norm \(2n\langle L_0\rangle+N(n^3-n)/12=O(N)\) by the commutator. At scale \(1/N\), either fluctuation has vanishing matrix elements against \(B\Psi\). In the total-mode identity, the first \(k\) terms contribute \(O_B(N^{-1})\) because \(B L_n\) is bounded there. All unseen sites contribute equally by exchangeability. The preceding uniform integrability now passes this test to the limiting \((k+1)\)-site hierarchy. It follows that \[
 \int (\langle L_n\rangle_\sigma-\delta_{n2}e^s)
                    {\rm tr}(\sigma^{\otimes k} B)\,d\mu(\sigma)=0 .
\] The same for the adjoint follows either by conjugation or its fluctuation estimate. The bounded tensor-cylinder functions generate the state Borel sigma-field: matrix entries of a countable increasing family of finite occupation projectors separate density operators, and products are tested by larger \(k\). A monotone-class argument for the integrable signed measures therefore gives \(\langle L_n\rangle_\sigma=\delta_{n2}e^s\) for \(\mu\)-almost every \(\sigma\), simultaneously in the countable set of \(n>0\). This avoids assuming the macroscopic \(L_0\) concentration in order to derive it later. Also \(\langle L_0\rangle_\sigma\) is finite almost everywhere by energy integrability.

\section{Pure Gaussian and the exact algebraic Grunsky converse}\label{pure-gaussian-and-the-exact-algebraic-grunsky-converse}

The diagonal two-site support is contained in the symmetric swap sector. Thus \(0=\int(1-{\rm tr}\,\sigma^2)/2\,d\mu\) forces every directing state to be almost surely pure. Each \(\psi^{\otimes k}\) of its corresponding vector is supported in the diagonal vacuum module for every finite \(k\). In particular it is fixed by the action of \(O(2)\) on two copies. For one-copy Bargmann entire generating function \(F(z)\), evaluated first on each finite set of mode coordinates, reflection gives \(F(-z)=F(z)\), and the \(45^\circ\) rotation with the second variable zero gives \[
                 F(z)F(0)=F(z/\sqrt2)^2 .
\] If \(F(0)=0\) this implies the vector is zero by density of finite coherent vectors. Otherwise the Taylor expansion of the local analytic logarithm forces its only nonconstant homogeneous term to have degree two. Thus the vector has the normalized squeezed form \[
             \psi(A)\ \propto\
                 \exp\!\left(\tfrac12\sum_{m,n\ge1}A_{mn}b_m^\dagger b_n^\dagger\right)\Omega
\] with a symmetric Hilbert--Schmidt matrix \(A\) and \(\|A\|<1\). Hilbert--Schmidt follows by its two-particle projection, and strict contractivity follows by testing a largest singular direction of the compact complex symmetric \(A\) (the one-mode squeezed norm diverges at modulus one or above). These conditions also suffice for the displayed vector to be normal.

It is not true that an arbitrary one-current pure Gaussian is a vacuum diffeomorphism state. The extra \textbf{all tensor vacuum supports} supply exactly the following polynomial relations. Adopt the exterior convention \[
 A_{mn}=-\sqrt{mn}\,b_{mn},\qquad
       \log\frac{g(z)-g(w)}{z-w}
                         =-\sum_{m,n\ge1}b_{mn}z^{-m}w^{-n}.             \tag{G}
\] For a formal product \(U=\prod_\ell \exp(s_\ell L_{-k_\ell})\) with \(k_\ell\ge2\), the vacuum coefficient of \(U\Omega\) is one. The current identity \([L_{-k},J(z)]=(z^{1-k}\partial_z+(1-k)z^{-k})J(z)\) gives a formal exterior change of variable \(g(z)=z+O(z^{-1})\) for \(U^{-1}J(z)U=g'(z)J(g(z))\). Translation of \(g\) is irrelevant. Since \(\langle\Omega|U=\langle\Omega|\), its vacuum-to-vector two-current expectation minus the vacuum contraction is \[
 \frac{g'(z)g'(w)}{(g(z)-g(w))^2}-\frac1{(z-w)^2}.
\] Only double negative powers \(z^{-m-1}w^{-n-1}\) survive this subtraction (commute any negative current past the annihilation bra; its scalar contraction is the vacuum term). The mixed derivative of (G) therefore gives the exact Gaussian two-particle coefficients \(-\sqrt{mn}\,b_{mn}\) of the corresponding formal state. This uses formal power series with finite bras, not a large parameter analytic continuation.

Every formal coefficient \(b_{mn}\) in (G) is a finite polynomial in \(b_{j1}\): expand the logarithm of the divided difference and read its first column \(g(z)=z+\sum_j b_{j1} z^{-j}\). To enforce any one of these polynomial identities on \(A\), realize its left-hand side as a finite bra on \(k\) tensor factors, \(k\) at least its degree. A monomial in \(A_{mn}\) is realized by putting \(\langle\Omega|b_m b_n\) on a distinct factor for each occurrence and a vacuum bra on unused factors. For a formal diagonal negative-mode product on the \(k\)-copy vacuum the pairing is the polynomial evaluated on the corresponding \(g\), hence zero. Taking parameter derivatives yields zero against every ordered PBW vacuum descendant, irrespective of degree, so this finite bra is orthogonal to the entire diagonal cyclic vacuum subspace. Its pairing with \(\psi(A)^{\otimes k}\) must vanish. Dividing by its nonzero vacuum coefficient to the \(k\)-th power recovers the desired polynomial identity. This proves (G) for all \(m,n\); no converse to an unproved infinite-dimensional orbit assertion is used.

Now \(g(z)=z+\sum_j b_{j1} z^{-j}\) is genuinely analytic for \(|z|>1\) by Hilbert--Schmidt summability. The double series in (G) converges normally there in both variables by the same Cauchy--Schwarz bound. Exponentiate the formal equality near infinity, and extend it analytically on the exterior product. The divided quotient is holomorphic also on the diagonal. It equals a nowhere-zero exponential, proving directly that \(g\) is univalent (and \(g'\ne0\)) on the exterior. This avoids any extra assumption about a boundary extension or a Grunsky-norm characterization.

\section{Dilation gives literal diffeomorphism rays and controls the stress topology}\label{dilation-gives-literal-diffeomorphism-rays-and-controls-the-stress-topology}

For \(0<\rho<1\) let \(g_\rho(z)=\rho g(z/\rho)\). It is univalent past the unit circle, so its boundary is an analytic Jordan curve with a smooth analytic interior Riemann map \(h_\rho\) and an analytic welding \(\gamma_\rho=g_\rho^{-1}\circ h_\rho\). Its exterior Grunsky squeeze matrix is exactly \(A_{\rho,mn}=\rho^{m+n}A_{mn}\). To identify the \textbf{literal} unitary vacuum ray, let \(F_n(t)\) be the monic Faber polynomial from \(\log[(g_\rho(z)-t)/z]=-\sum_{n\ge1}F_n(t)z^{-n}/n\); putting \(t=g_\rho(w)\) gives \(F_n(g_\rho(w))=w^n+n\sum_m b_{\rho,nm}w^{-m}\). Under current covariance the exterior smearing by this function is transported via welding to \(F_n(h_\rho(\cdot))\), analytic on the interior disk; its modes annihilate the vacuum. Thus the welded vacuum at the outer circle is annihilated by \[
                  a_n+n\sum_m b_{\rho,nm}a_{-m}.
\] Here \(a_n=(2\pi i)^{-1}\oint z^n J(z)dz\), and the zero-charge \(a_0=0\). The normal squeeze \(A_\rho=-\sqrt{mn}\,b_{\rho,mn}\) obeys exactly the same equations. They determine the state up to scalar by iterating its particle sectors from the vacuum coefficient. Therefore \(\psi(A_\rho)\) is precisely a smooth projective vacuum diffeomorphism ray, taking the inverse in the opposite covariance convention. This is also consistent with the exterior symplectic block and metaplectic welding construction of Alekseev--Shatashvili--Takhtajan \cite[Proposition 3.3, equation (30), Theorem 5.5]{welding}. The new step here is forcing the exterior block by finite tensor support, not that Gaussian-welding formula itself.

For stress convergence put \(D=\operatorname{diag}(1,2,\ldots)\) on normalized oscillators. The normal covariance \(K\) has weighted trace \({\rm tr}(D^{1/2}KD^{1/2})=E<\infty\). In one consistent index convention \[
 K=A^\dagger(1-AA^\dagger)^{-1}A,\qquad
 M=A(1-A^\dagger A)^{-1}=(1-AA^\dagger)^{-1}A
\] for the anomalous covariance. Thus \(AD^{1/2}\) is Hilbert--Schmidt by positivity of the resolvent, and the other weighted direction follows by symmetry. Dilation is \(A_\rho=D_\rho A D_\rho\) with \(D_\rho=\operatorname{diag}(\rho^n)\) and keeps the same strict norm bound. It converges in both energy-weighted Hilbert--Schmidt norms (bounded strong convergence of \(D_\rho\) on either side of a Hilbert--Schmidt operator). Resolvents converge in ordinary operator norm. In particular \(D^{1/2}K_\rho D^{1/2}\to D^{1/2}KD^{1/2}\) in trace norm by the positive factorization above, and \(M_\rho\to M\) in both weighted Hilbert--Schmidt norms by the two resolvent formulas. Transposes or complex conjugates from labeling the contractions have no effect on these bounds.

The normal contribution to the stress Fourier mode \(n\ge1\) is \(\sum_m Q_{m,m+n}\) with \(Q=D^{1/2}KD^{1/2}\). The difference between two such contributions is bounded uniformly in \(n\) by the trace norm of the difference of \(Q\), since each is a trace against a partial shift of norm one. The anomalous contribution at that mode is proportional to \(\sum_{m=1}^{n-1}\sqrt{m(n-m)}\,M_{m,n-m}\). Its difference squared divided by \(n^2\) is bounded by \(C\sum_{m<n}n|\Delta M_{m,n-m}|^2\), by Cauchy--Schwarz and \(\sum_{m<n}m(n-m)/n=O(n^2)\). Sum over all \(n\) using the weighted norms. The mean mode converges by the weighted trace norm. Thus the cylinder stress expectations of \(\psi(A_\rho)\) converge strongly in real circle \(H^{-1}\) to those of \(\psi(A)\).

The antiperiodic Hill quadratic forms are continuous under real \(H^{-1}\) perturbations, including their first two eigenvalues. One direct argument: for \(u\in H^1\) antiperiodic, \(|u|^2\) is periodic and \(\||u|^2\|_{H^1}\le C\|u\|_{H^1}^2\), so \(H^{-1}\) differences are form-small as \(\rho\uparrow1\). For a fixed \(H^{-1}\) reference, use its \(L^2\) zero-mean primitive and one-dimensional interpolation to obtain infinitesimal form-boundedness; this supplies a common lower bound and equivalence of the positive shift form with the \(H^1\) norm. The compact min--max principle gives the claimed eigenvalue convergence. For a circle diffeomorphism with increasing lift \(F\), the normalized vacuum-orbit potential \(q_F=6t_F=-\tfrac14(F')^2-\tfrac12\{F,\theta\}\) satisfies the Liouville form identity
\[
 \int_0^{2\pi} (|u'|^2+q_F|u|^2)\,d\theta
 =\int_0^{2\pi} (|v'|^2-\tfrac14|v|^2)\,dx,\qquad
 u=(F')^{-1/2}v\circ F .
\]
It follows by the chain rule and periodic integration by parts; the transformation is an isomorphism of the antiperiodic \(H^1\) form domains. Thus the first two antiperiodic eigenvalues of \(-\partial_\theta^2+6t_{A_\rho}\) are exactly zero. The limit stress has only the fixed quadrupole and its mean, by the Ward argument, so it is smooth and retains the double zero. In detail, \(-\partial_\theta^2+12\alpha\cos(2\theta)\) on antiperiodic functions on the whole \(2\pi\) circle is the direct sum of the \(\pm i\) conditions on a \(\pi\)-cell, since translation by \(\pi\) squares to \(-I\). The ground eigenvalues of the two blocks coincide by conjugation. Each is simple: a real determinant-one monodromy having \(i\) also has the distinct multiplier \(-i\), with one eigenline for \(i\). Identifying the zero of the antiperiodic bottom therefore fixes the normalized mean to be \[
            e_*(\alpha)=\tfrac16\bigl(\tfrac14-a_*(6\alpha)\bigr),
            \qquad \alpha=e^s ,
\] where \(a_*(q)\) is the lowest \(i\)-quasiperiodic Mathieu eigenvalue on a \(\pi\) cell. No smoothness of the undilated weld, or uniform trace-norm compactness of diffeomorphism group elements, was assumed or needed.

\section{\texorpdfstring{Equality at all real \(c\); true stress concentration}{Equality at all real c; true stress concentration}}\label{equality-at-all-real-c-true-stress-concentration}

Every directing component of every integer-\(N\) subsequential hierarchy has the same one-site energy \(e_*(\alpha)\). The two no-escape estimates therefore give
\[
 \lim_{N\to\infty}\frac{\langle L_0\rangle_{N,N\alpha}}N=e_*(\alpha)
 \qquad(\alpha>0).
\] The continuous Fekete limit is shared with that integer subsequence. Convexity bounds \(f_N'\) uniformly on each compact \(s\)-interval using \(0\le f_N\le2e^s\). Hence the fundamental theorem of calculus and dominated convergence pass the integer energy limit to \(f(s)-f(s_0)\); let \(s_0\to-\infty\) using the same bound. This proves \[
 f(s)=\int_{-\infty}^s e_*(e^u)\,du,\qquad
 \lim_{c\to\infty,\ c\in\mathbb R}\frac{\langle L_0\rangle_{c,c\alpha}}c
                      = e_*(\alpha).
\] The real-charge derivative convergence follows from convexity and differentiability of the limiting \(f\); analyticity of \(a_*\) on the real line follows from its simple twisted eigenvalue. Thus the normalized gap is zero on the entire positive macroscopic ray. This does not contradict the strict finite-charge exclusion in \cite{p34}. We give its short orbit comparison below; no finite correction to the gap is asserted.

For the actual stress concentration use the exact tilt \(\mathbb E e^{tL_0}=\exp(c[f_c(s+t)-f_c(s)])\). Differentiability of \(f\) yields the two Chernoff estimates for \(L_0/c\) about \(e_*\), with exponential large-\(c\) tails and convergence of variance by the moment bounds in Section 2. For positive \(n\) the annihilator fluctuation on the centered ket vanishes, and the adjoint fluctuation has norm squared \(2n\langle L_0\rangle+c(n^3-n)/12\). For a fixed smooth real smearing the weighted sum of these norms is \(O(\sqrt c)\) because its Fourier coefficients decay rapidly; the zero-mode part is \(o(c)\) in fluctuation norm. Therefore the normalized stress smearing variance tends to zero for every fixed smooth test; it is stronger than numerical agreement of stress means.

\section{Smooth classical representative and strict finite-charge comparison}
\label{sec:classical}

The lowest twisted eigenfunction \(y\) on the \(\pi\)-cell at \(q=6\alpha\) is simple. Its Wronskian \(\operatorname{Im}(\bar y y')\) is nonzero, because its conjugate has the distinct multiplier \(-i\). It is positive by continuous continuation of the lowest simple eigenfunction from \(q=0\), where \(y=e^{i\theta/2}\). As a result \(y\) has no zeros, and its unwrapped phase increment over a cell is \(\pi/2\) (the integer multiple of \(2\pi\) in that increment is constant along the continuation). Therefore \(F=2\arg y\) is an increasing degree-one lift. Constancy of the Wronskian gives \(|y|\propto(F')^{-1/2}\); direct differentiation gives
\[
 \frac{y''}{y}
 =-\tfrac14(F')^2-\tfrac12\{F,\theta\}
 =-a_*(6\alpha)+12\alpha\cos(2\theta).
\]
This is the exact Schwarzian transformation of the normalized vacuum cylinder stress by \(F\) \cite{fewster}. It constructs the smooth ray \(P_\alpha\) in Corollary~\ref{cor:product}.

For any state with the same quadrupole and energy \(E\), its inverse diffeomorphism changes its stress to the vacuum value plus \((E-c e_*(\alpha))\) times the positive squared inverse derivative. This follows by subtracting the two stress transformation laws so that their Schwarzian terms cancel. Its rotation energy is nonnegative. Equality is possible only on the unique vacuum ray after the inverse diffeomorphism. This also proves uniqueness in the \(c=1\) current space, whose zero-momentum rotation operator has simple eigenvalue zero. This transfer may be made at the level of quadratic forms, without assuming a common operator domain for the limiting vector. For a smooth stress smearing \(w\), the normal current term is bounded by a constant times \(L_0\) as a form: on one-particle space its energy-conjugated kernel is the bounded Toeplitz matrix of \(w\). The anomalous current contribution is a sum over finitely many pairs for each Fourier index, and
\[
 |\langle a_j a_k\rangle|
 \le \big[\langle a_j a_{-j}\rangle\langle a_{-k}a_k\rangle\big]^{1/2}
 \le C(E+j+k)
\]
on a normalized vector of energy \(E\). Rapid decay of the Fourier coefficients makes the sum an \(O(1+E)\) form, with polarization extending it continuously to the energy form domain. Smooth covariance for finite-energy vectors then implies that \(U(F)^{\pm1}\) is bounded on this form domain, since \(U(F)L_0U(F)^{-1}\) is a smooth stress smearing plus a scalar. Approximation in the energy form norm gives the claimed covariance for every directing state. Thus all directing components equal \(P_\alpha\), so every trace-norm subsequential \(k\)-copy limit is \(P_\alpha^{\otimes k}\), proving the corollary.

For clarity, here is the exclusion of equality for the finite-charge Whittaker vector, also proved in \cite{p34}. If \(U(F)\Omega\) were a simultaneous positive-mode eigenvector, its pulled-back vector fields \(g_n=e^{inF}/F'\) would generate generators mapping the vacuum to a scalar multiple of itself. The modes \(L_{-j}\Omega\) at \(j\ge2\) have nonzero norms and different energies, so the Fourier coefficients of \(g_n\) below \(-1\) vanish. Each \(G_n(e^{i\theta})=e^{i\theta}g_n(\theta)\) therefore extends holomorphically to the disk, continuously and smoothly to its boundary. The identities \(G_nG_1^{n-2}=G_2^{n-1}\) imply \(G_2/G_1\) has no interior poles: a pole of fixed order cannot be canceled by \(G_1\) in \(G_n=G_1(G_2/G_1)^{n-1}\) for all \(n\). On the boundary this quotient equals \(e^{iF}\), of modulus one and winding one. By the argument principle and the maximum principle it is a single Blaschke factor times a phase. The corresponding circle Möbius transformation preserves the vacuum ray. This contradicts \(r>0\) for \(\Psi_{c,r}\), giving strictness.

\begin{thebibliography}{99}
\bibitem{p34} \emph{Whittaker states outside the classical Virasoro vacuum orbit}, \href{https://github.com/g3173028019-beep/ai-automated-quantum-gravity-research/tree/main/papers/34-ads3-vacuum-stress}{source and verification}.
\bibitem{bh} J. D. Brown and M. Henneaux, \emph{Central charges in the canonical realization of asymptotic symmetries}, \href{https://doi.org/10.1007/BF01211590}{Commun. Math. Phys. 104 (1986), 207--226}.
\bibitem{cotler} J. Cotler and K. Jensen, \emph{A theory of reparameterizations for AdS$_3$ gravity}, \href{https://arxiv.org/abs/1808.03263}{arXiv:1808.03263}.
\bibitem{definetti} M. Lewin, P. T. Nam and N. Rougerie, \emph{Remarks on the quantum de Finetti theorem for bosonic systems}, \href{https://arxiv.org/abs/1310.2200}{arXiv:1310.2200}.
\bibitem{voa} S. Carpi, Y. Kawahigashi, R. Longo and M. Weiner, \emph{From vertex operator algebras to conformal nets and back}, \href{https://arxiv.org/abs/1503.01260}{arXiv:1503.01260}.
\bibitem{fewster} C. J. Fewster and S. Hollands, \emph{Quantum energy inequalities in two-dimensional conformal field theory}, \href{https://arxiv.org/abs/math-ph/0412028}{arXiv:math-ph/0412028}.
\bibitem{carpi} S. Carpi, S. Del Vecchio, S. Iovieno and Y. Tanimoto, \emph{Positive energy representations of Sobolev diffeomorphism groups of the circle}, \href{https://arxiv.org/abs/1808.02384}{arXiv:1808.02384}.
\bibitem{welding} A. Alekseev, S. Shatashvili and L. Takhtajan, \emph{Berezin quantization, conformal welding and the Bott--Virasoro group}, \href{https://arxiv.org/abs/2212.04822}{arXiv:2212.04822}.
\bibitem{piatek14} M. Pi\k atek and A. R. Pietrykowski, \emph{Classical irregular block, \(\mathcal N=2\) pure gauge theory and Mathieu equation}, \href{https://arxiv.org/abs/1407.0305}{arXiv:1407.0305}.
\bibitem{piatek15} M. Pi\k atek and A. R. Pietrykowski, \emph{Classical limit of irregular blocks and Mathieu functions}, \href{https://arxiv.org/abs/1509.08164}{arXiv:1509.08164}.
\bibitem{nekrasov} P. Arnaudo, G. Bonelli and A. Tanzini, \emph{On the convergence of Nekrasov functions}, \href{https://arxiv.org/abs/2212.06741}{arXiv:2212.06741}.
\end{thebibliography}
\end{document}
